% Créditos de las fotografías usadas en este volumen (Libro 1).
% Los registros de licencia legibles por máquina están en images/book1/CREDITS.md;
% mantener ambos archivos sincronizados al añadir o quitar una fotografía.
\chapter*{Créditos de las imágenes}

Los dibujos de este libro son originales. Las fotografías se usan bajo
sus respectivas licencias libres, con nuestro agradecimiento a sus
autores:

\begin{itemize}
  \item \emph{John Dalton} (capítulo de los átomos y las moléculas):
        grabado de Charles Turner a partir de Joseph Allen, 1834, Library
        of Congress, vía Wikimedia Commons, dominio público.
  \item \emph{La tabla de los elementos de Dalton, 1808} (capítulo de los
        átomos y las moléculas): lámina de \emph{A New System of Chemical
        Philosophy} de John Dalton, digitalización del Internet Archive,
        vía Wikimedia Commons, dominio público.
  \item \emph{Bauxita} (capítulo de las materias primas): fotografía de
        Werner Schellmann, vía Wikimedia Commons, CC BY-SA 2.5.
  \item \emph{Granito} (capítulo de las sustancias puras): fotografía de
        Eurico Zimbres, vía Wikimedia Commons, CC BY-SA 2.0 br.
  \item \emph{Sulfato de cobre anhidro} (capítulo de la identificación de
        sustancias): fotografía de W.~Oelen, vía Wikimedia Commons,
        CC BY-SA 3.0.
  \item \emph{Cristales de sulfato de cobre} (capítulo de la
        identificación de sustancias): fotografía de Stephanb, vía
        Wikimedia Commons, CC BY-SA 3.0.
  \item \emph{Antoine-Laurent Lavoisier y Marie-Anne Paulze Lavoisier}
        (capítulo de las ecuaciones ajustadas): cuadro de Jacques-Louis
        David, 1788, Metropolitan Museum of Art, vía Wikimedia Commons,
        dominio público.
  \item \emph{Teléfono de baquelita} (capítulo de los plásticos):
        fotografía de Holger Ellgaard, vía Wikimedia Commons,
        CC BY-SA 3.0.
  \item \emph{Ernest Rutherford} (capítulo del interior del átomo):
        fotografía, George Grantham Bain Collection, Library of Congress,
        vía Wikimedia Commons, dominio público.
  \item \emph{Halita} (capítulo de los iones): fotografía de Robert
        M.~Lavinsky, vía Wikimedia Commons, CC BY-SA 3.0.
  \item \emph{La Estatua de la Libertad} (capítulo de los metales y la
        corrosión): fotografía de Elcobbola, vía Wikimedia Commons,
        dominio público.
  \item \emph{Dmitri Mendeléiev} y \emph{su tabla de 1869} (capítulo de la
        tabla periódica): vía Wikimedia Commons, dominio público.
  \item \emph{La nebulosa del Cangrejo} (capítulo de los elementos): NASA,
        ESA, J.~Hester y A.~Loll, dominio público.
  \item \emph{Sodio} (capítulo de las capas electrónicas): fotografía de
        Dnn87, vía Wikimedia Commons, CC BY-SA 3.0.
  \item \emph{Cloro en una ampolla} (capítulo de las capas electrónicas):
        fotografía de W.~Oelen, vía Wikimedia Commons, CC BY-SA 3.0.
  \item \emph{Amedeo Avogadro} (capítulo del mol): litografía a partir de
        un dibujo de C.~Sentier, 1856, colección Edgar Fahs Smith, vía
        Wikimedia Commons, dominio público.
  \item \emph{Cristales de nieve} (capítulo de la polaridad): fotografías
        de Wilson Bentley, 1902, vía Wikimedia Commons, dominio público.
  \item \emph{Louis Pasteur} (capítulo de estereoquímica): fotografía de
        Paul Nadar, vía Wikimedia Commons, dominio público.
  \item \emph{Cristales de tartrato de amonio} (capítulo de
        estereoquímica): fotografía de Marie-Lan Taÿ Pamart, vía Wikimedia
        Commons, CC BY 4.0.
  \item \emph{Svante Arrhenius} (capítulo de los ácidos y las bases):
        fotograbado, 1909, vía Wikimedia Commons, dominio público.
  \item \emph{Una pila voltaica} (capítulo de las pilas): fotografía de
        GuidoB, vía Wikimedia Commons, CC BY-SA 3.0.
  \item \emph{Wallace Carothers en su laboratorio} (capítulo de los
        polímeros): fotógrafo desconocido, vía Wikimedia Commons, dominio
        público.
\end{itemize}

Los enlaces completos a las fuentes y las URL de las licencias de cada
fotografía figuran en \texttt{images/book1/CREDITS.md}, en el repositorio
del libro. Los pictogramas de peligro son los del Sistema Globalmente
Armonizado de Clasificación y Etiquetado de Productos Químicos, tal como
los dibuja el paquete \texttt{ghsystem} de \TeX\ Live.

\bigskip
Junto a las fotografías y los dibujos originales, algunas figuras son
ilustraciones generadas por IA, producidas con la generación de imágenes
de OpenAI a partir de instrucciones escritas por los editores del libro,
y revisadas desde el punto de vista químico antes de incluirlas. La
instrucción completa de cada imagen generada está registrada en
\texttt{images/book1/ai/PROMPTS.md}, en el repositorio.
\clearpage
